%===============================================================================
% LLMWiki Technical Specification — Volume 03: Knowledge Graph
%===============================================================================
\chapter{Knowledge Graph \& Semantic Index}
\label{cha:knowledge-graph}

This volume specifies the \emph{knowledge graph} (property graph with typed nodes/edges) and the \emph{semantic index} (hybrid search: vector + keyword + graph adjacency). These two modules form Layer 2 Core (Vol.~01, \S\ref{sec:arch:core}) and are tightly coupled: the graph provides structural semantics, the index provides fast retrieval, and both share the same content-addressed unit identifiers.

\section{Knowledge Graph}
\label{sec:kg:graph}

\subsection{Data Model}
\label{sec:kg:data-model}

\begin{definition}[Property Graph]
\label{def:property-graph}
A \emph{property graph} $G = (V, E, \lambda_V, \lambda_E, \sigma_V, \sigma_E)$ where:
\begin{itemize}
  \item $V$: set of nodes (entities, concepts, units)
  \item $E \subseteq V \times V$: set of directed edges (relationships)
  \item $\lambda_V: V \to \mathcal{L}_V$: node label function (type)
  \item $\lambda_E: E \to \mathcal{L}_E$: edge label function (relation type)
  \item $\sigma_V: V \to \mathcal{P}(\text{Key} \times \text{Value})$: node properties
  \item $\sigma_E: E \to \mathcal{P}(\text{Key} \times \text{Value})$: edge properties
\end{itemize}
\end{definition}

\paragraph{Node Types.}
\begin{itemize}
  \item \texttt{Unit}: A semantic unit from parsing (Vol.~02). Properties: \texttt{id (SHA256)}, \texttt{type}, \texttt{span}, \texttt{file\_path}, \texttt{language}, \texttt{metadata(JSON)}.
  \item \texttt{Entity}: Resolved entity (class, function, table, concept). Properties: \texttt{name}, \texttt{type}, \texttt{canonical\_id}, \texttt{description}, \texttt{embedding\_ref}.
  \item \texttt{File}: Source artifact. Properties: \texttt{path}, \texttt{hash}, \texttt{size}, \texttt{mtime}, \texttt{mime}.
  \item \texttt{Module/Package/Namespace}: Logical grouping. Properties: \texttt{name}, \texttt{parent\_id}.
\end{itemize}

\paragraph{Edge Types.}
\begin{itemize}
  \item \texttt{CONTAINS}: File $\to$ Unit, Module $\to$ Unit, Unit $\to$ Unit (nesting)
  \item \texttt{DEFINES}: Unit $\to$ Entity (unit defines entity)
  \item \texttt{REFERENCES}: Unit $\to$ Entity (unit references entity)
  \item \texttt{CALLS}: Entity $\to$ Entity (function call)
  \item \texttt{INHERITS}: Entity $\to$ Entity (class inheritance)
  \item \texttt{IMPORTS}: Module $\to$ Module
  \item \texttt{RELATES}: Entity $\leftrightarrow$ Entity (semantic similarity, co-occurrence)
  \item \texttt{NEXT\_TO}: Unit $\to$ Unit (sequential proximity in source)
\end{itemize}

\paragraph{Content-Addressed Identity.}
Every \texttt{Unit} node has \texttt{id = SHA256(canonical\_form)}. \texttt{Entity} nodes have \texttt{canonical\_id} derived from name + type + context hash. This ensures deterministic identity across re-ingestion.

\subsection{Graph Store Backend}
\label{sec:kg:store}

The \texttt{GraphStore} trait (Vol.~01, \S\ref{iface:graphstore}) admits multiple backends:

\begin{table}[htbp]
\centering
\caption{Graph Store Backend Comparison}
\label{tab:kg:backends}
\begin{adjustbox}{max width=\textwidth}
\begin{tabular}{lcccc}
\toprule
\textbf{Backend} & \textbf{Embedded} & \textbf{Scale} & \textbf{Transactions} & \textbf{License} \\
\midrule
RocksDB (custom) & Yes & 10B+ edges & Yes (WAL) & Apache-2.0 \\
SQLite (with ext) & Yes & 100M edges & Yes (WAL) & Public Domain \\
Redwood (Rust) & Yes & 1B edges & Yes & MIT \\
Neo4j Embedded & No & 100B edges & Yes & GPL/Commercial \\
Kuzu & Yes & 10B edges & Yes & MIT \\
\bottomrule
\end{tabular}
\end{adjustbox}
\end{table}

\textbf{Default:} RocksDB-based custom store (\texttt{llmwiki-graph}). Nodes and edges stored as separate column families with prefix encoding for efficient prefix scans (by type, by partition).

\subsubsection{Storage Layout (RocksDB)}
\label{sec:kg:rocksdb-layout}

\begin{verbatim}
Column Families:
|-- nodes:{type}       # Key: node_id (fixed 32 bytes), Value: serialized props
|-- edges:{rel_type}   # Key: src_id(32) | dst_id(32), Value: edge props
|-- adj:out:{type}     # Key: src_id, Value: list of (rel_type, dst_id) - adjacency
|-- adj:in:{type}      # Key: dst_id, Value: list of (rel_type, src_id) - reverse adj
|-- idx:label:{label}  # Key: label | node_id, Value: empty (label index)
|-- idx:prop:{key}     # Key: prop_value | node_id, Value: empty (property index)
`-- meta               # Key: "stats", "version", "config", Value: JSON
\end{verbatim}

\subsection{Entity Resolution}
\label{sec:kg:entity-resolution}

When a parser emits a \texttt{Unit} that \texttt{DEFINES} or \texttt{REFERENCES} an entity, the graph layer resolves it to a canonical \texttt{Entity} node.

\begin{algorithm}[H]
\caption{Entity Resolution}
\label{alg:entity-resolution}
\begin{algorithmic}[1]
\Require Unit $u$ with entity mention $m = (name, type, context)$
\Ensure Resolved \texttt{Entity} node $e$
\State $candidates \gets \textsc{LookupByNameType}(G, m.name, m.type)$
\If{$candidates = \varnothing$}
  \State \Return \textsc{CreateEntity}(m, u)
\EndIf
\State $best \gets \arg\max_{c \in candidates} \textsc{Score}(c, m, u)$
\If{$\textsc{Score}(best, m, u) \geq \tau_{er}$}
  \State \textsc{UpdateEntity}(best, u) \Comment{Merge context, update embedding}
  \State \Return $best$
\Else
  \State \Return \textsc{CreateEntity}(m, u) \Comment{New distinct entity}
\EndIf
\end{algorithmic}
\end{algorithm}

\textbf{Scoring function:}
\begin{align*}
\textsc{Score}(c, m, u) = &\; \alpha \cdot \text{FuzzyMatch}(c.name, m.name) \\
&+ \beta \cdot \text{TypeCompat}(c.type, m.type) \\
&+ \gamma \cdot \text{ContextOverlap}(c.context, u.context) \\
&+ \delta \cdot \text{CosineSim}(c.embedding, u.embedding)
\end{align*}
Default weights: $\alpha=0.3, \beta=0.2, \gamma=0.2, \delta=0.3$. Threshold $\tau_{er}=0.85$ (configurable).

\subsection{Incremental Graph Updates}
\label{sec:kg:incremental}

On ingestion diff $(added, modified, removed)$ units (Vol.~02, Alg.~\ref{alg:incremental-diff}):

\begin{algorithm}[H]
\caption{Incremental Graph Update}
\label{alg:graph-incremental}
\begin{algorithmic}[1]
\Require Diff sets $A$ (added), $M$ (modified), $R$ (removed)
\Ensure Updated graph $G'$
\ForEach{$u \in R$}
  \State \textsc{DeleteNode}(u.id) \Comment{Tombstone}
  \State \textsc{DeleteEdges}(u.id) \Comment{All incident edges}
\EndFor
\ForEach{$u \in A \cup M$}
  \State \textsc{UpsertUnitNode}(u)
  \State $entities \gets \textsc{ExtractEntities}(u)$
  \ForEach{$e \in entities$}
    \State $resolved \gets \textsc{ResolveEntity}(e, u)$ \Comment{Alg.~\ref{alg:entity-resolution}}
    \If{$e.relation = \text{DEFINES}$}
      \State \textsc{UpsertEdge}(u.id, resolved, \texttt{DEFINES})
    \Else
      \State \textsc{UpsertEdge}(u.id, resolved, \texttt{REFERENCES})
    \EndIf
  \EndFor
  \State \textsc{UpsertAdjacency}(u) \Comment{Contains, Next\_To}
\EndFor
\State \textsc{CommitTransaction}()
\State \Return $G$
\end{algorithmic}
\end{algorithm}

All operations within a single ingestion are batched in one RocksDB write batch (atomic).

\subsection{Graph Traversal Primitives}
\label{sec:kg:traversal}

The \texttt{traverse} and \texttt{neighborhoods} methods (Vol.~01, \S\ref{iface:graphstore}) support:

\begin{itemize}
  \item \textbf{Pattern matching}: Cypher-like declarative patterns (e.g., \texttt{(a:Function)-[:CALLS*1..3]->(b:Function)})
  \item \textbf{Bounded BFS/DFS}: $k$-hop neighborhoods with edge-type filters
  \item \textbf{Personalized PageRank}: For HippoRAG-style retrieval (Vol.~04)
  \item \textbf{Shortest path}: Between two entities
  \item \textbf{Subgraph extraction}: Induced subgraph for context packing
\end{itemize}

\begin{verbatim}
# TraversalPattern examples
pattern = {
  "start": { "type": "Function", "name": "process_order" },
  "hops": 3,
  "edge_types": ["CALLS", "REFERENCES"],
  "direction": "out",
  "filter": { "node_type": ["Function", "Class"] },
  "limit": 100
}
\end{verbatim}

\section{Semantic Index}
\label{sec:kg:index}

The \texttt{SemanticIndex} (Vol.~01, \S\ref{iface:semantic-index}) is a hybrid index combining three retrieval signals:

\begin{enumerate}
  \item \textbf{Vector (HNSW)}: Dense embedding similarity
  \item \textbf{Keyword (BM25)}: Sparse lexical matching
  \item \textbf{Graph Adjacency}: Structural proximity in knowledge graph
\end{enumerate}

\subsection{HNSW Vector Index}
\label{sec:kg:hnsw}

\paragraph{Implementation:} \texttt{hnswlib} (C++ with Python/Rust bindings) or \texttt{faiss} (with disk-backed index). Disk-backed mode: memory-mapped file for vectors, metadata in separate store.

\paragraph{Configuration:}
\begin{verbatim}
hnsw:
  M: 16                    # max connections per node
  ef_construction: 200     # build-time search width
  ef_search: 64           # query-time search width
  dim: 1024               # embedding dimension
  space: "cosine"         # or "l2", "ip"
  disk_backed: true
  mmap: true
\end{verbatim}

\paragraph{Content-Addressed Vectors.}
Vectors are stored keyed by \texttt{unit\_id} (SHA256). The index maps \texttt{unit\_id} $\to$ internal HNSW label. On upsert: if unit exists, update vector; if new, add. On delete: mark tombstone (lazy cleanup during compaction).

\subsection{BM25 Keyword Index}
\label{sec:kg:bm25}

\paragraph{Implementation:} \texttt{tantivy} (Rust) or \texttt{OpenSearch} (if distributed). Tantivy is embedded, fast, and supports incremental updates.

\paragraph{Schema:}
\begin{verbatim}
fields:
  - unit_id: Keyword (stored, indexed)  # SHA256
  - content: Text (analyzer: language-specific)
  - file_path: Keyword (stored, indexed)
  - unit_type: Keyword (stored, indexed)
  - language: Keyword (stored, indexed)
\end{verbatim}

\paragraph{Analyzer:} Language-specific (e.g., \texttt{code} analyzer for source: split on punctuation, camelCase, snake\_case; \texttt{standard} for prose).

\subsection{Graph Adjacency Index}
\label{sec:kg:adjacency-index}

For query-time graph expansion (Vol.~04), we maintain a lightweight adjacency index separate from the full graph store:

\begin{verbatim}
Adjacency Index (RocksDB column family "adj:index"):
Key: unit_id (32 bytes)
Value: serialized adjacency list
  [
    { "rel": "CONTAINS", "target": "sha256...", "weight": 1.0 },
    { "rel": "REFERENCES", "target": "sha256...", "weight": 0.8 },
    { "rel": "NEXT_TO", "target": "sha256...", "weight": 0.3 }
  ]
\end{verbatim}

Updated incrementally alongside graph. Used for \texttt{search\_graph} in hybrid retrieval.

\subsection{Hybrid Search}
\label{sec:kg:hybrid-search}

The \texttt{search\_hybrid} method combines three signals:

\begin{algorithm}[H]
\caption{Hybrid Search}
\label{alg:hybrid-search}
\begin{algorithmic}[1]
\Require Query $q$ (text + optional embedding), $k$, weights $w_v, w_k, w_g$
\Ensure Top-$k$ scored units
\State $q_v \gets \textsc{Embed}(q.text)$ \Comment{If not provided}
\State $R_v \gets \textsc{HNSWSearch}(q_v, k \times 3)$
\State $R_k \gets \textsc{BM25Search}(q.text, k \times 3)$
\State $R_g \gets \textsc{GraphExpand}(q, k \times 2)$ \Comment{Adjacency traversal}
\State $scores \gets \varnothing$
\ForEach{$r \in R_v \cup R_k \cup R_g$}
  \State $s_v \gets r \in R_v ? \text{Score}(r, R_v) : 0$
  \State $s_k \gets r \in R_k ? \text{Score}(r, R_k) : 0$
  \State $s_g \gets r \in R_g ? \text{Score}(r, R_g) : 0$
  \State $scores[r.id] = w_v \cdot s_v + w_k \cdot s_k + w_g \cdot s_g$
\EndFor
\State \Return \textsc{TopK}($scores, k$)
\end{algorithmic}
\end{algorithm}

Default weights: $w_v=0.6, w_k=0.3, w_g=0.1$ (configurable per query type).

Score normalization: Reciprocal Rank Fusion (RRF) or min-max normalization per signal before weighted sum.

\subsection{Index Update Strategy}
\label{sec:kg:index-update}

\paragraph{Vector (HNSW):} Incremental add/update. HNSW supports dynamic insertion. Deletes = tombstones. Periodic \texttt{rebuild\_threshold} triggers full rebuild.

\paragraph{BM25 (Tantivy):} Incremental commit. Segments merged in background. Delete = term tombstone.

\paragraph{Graph Adjacency:} Updated in same transaction as graph (RocksDB batch).

\section{Query-Time Graph Expansion}
\label{sec:kg:graph-expansion}

Used by retriever to expand candidate set via graph structure:

\begin{enumerate}
  \item \textbf{Seed}: Start from vector/keyword hits $S$
  \item \textbf{Expand}: For each $s \in S$, fetch $k$-hop neighborhood $N_k(s)$
  \item \textbf{Score}: Personalized PageRank from $S$ over induced subgraph
  \item \textbf{Merge}: Add top-$m$ PPR nodes to candidate pool
\end{enumerate}

This implements the ``graph traversal'' retrieval path complementary to vector/keyword.

\section{Performance Characteristics}
\label{sec:kg:performance}

\begin{table}[htbp]
\centering
\caption{Expected Performance (Single Node, SSD)}
\label{tab:kg:perf}
\begin{adjustbox}{max width=\textwidth}
\begin{tabular}{lcc}
\toprule
\textbf{Operation} & \textbf{Latency (p99)} & \textbf{Throughput} \\
\midrule
Node upsert (batch 1000) & < 50 ms & 20k nodes/s \\
Edge upsert (batch 5000) & < 100 ms & 50k edges/s \\
$k$-hop traversal (k=2, 100 seeds) & < 10 ms & 10k queries/s \\
PPR (100 seeds, 3 hops) & < 50 ms & 2k queries/s \\
HNSW search (k=10, ef=64) & < 5 ms & 2k queries/s \\
BM25 search (k=10) & < 3 ms & 3k queries/s \\
Hybrid search (k=20) & < 15 ms & 1k queries/s \\
\bottomrule
\end{tabular}
\end{adjustbox}
\end{table}

\section{Patent-Relevant Novelty}
\label{sec:kg:patent}

\begin{patentclaim}
\textbf{Claim 4:} A knowledge graph system wherein: (a) all nodes are content-addressed via SHA-256 of their canonical representation, (b) entity resolution merges mentions using a weighted combination of name similarity, type compatibility, context overlap, and embedding similarity, (c) the graph is persisted in a column-family store with prefix-encoded keys enabling type-partitioned scans, and (d) a separate graph adjacency index enables sub-millisecond $k$-hop expansion for hybrid retrieval.
\end{patentclaim}

\begin{patentclaim}
\textbf{Claim 5:} The system of Claim 4 wherein incremental updates to the graph and indices are executed atomically within a single write batch, and deterministic replay is achieved by re-applying the ingestion log against an empty store.
\end{patentclaim}

\section{Summary}
\label{sec:kg:summary}

The knowledge graph and semantic index form the queryable substrate of LLMWiki:
\begin{enumerate}
  \item \textbf{Graph}: Property graph with content-addressed nodes, entity resolution, and incremental updates
  \item \textbf{Vector Index}: HNSW with disk-backed storage, SHA-256 keyed
  \item \textbf{Keyword Index}: Tantivy BM25 with language-aware analyzers
  \item \textbf{Adjacency Index}: Lightweight graph expansion for structural retrieval
  \item \textbf{Hybrid Search}: RRF-weighted fusion of three signals
\end{enumerate}
All indices share the same unit identifiers, enabling seamless cross-index joins and evidence tracing (Vol.~06).

The next volume (Vol.~04) details the retrieval pipeline including the Evidence Gate.

%===============================================================================
% End of Volume 03
%===============================================================================
